\metatitle{Simple vertices, stability and matching polynomials}
\metadescription{One simple assignment vertex gives stable weighted completion
polynomials, a matching representation, and real-rooted counting polynomials.}

\section[simpleVertexStability]{From one simple vertex to real stability}

The geometry of one assignment vertex can control every weighted color
completion. The resulting representation connects colored matroids with
bipartite matchings and the classical real-rootedness of rook polynomials.

Let $\mathcal F$ be a colored matroid on $E$, fix a target color $c$, and
give every hidden cell $(e,a)$, $a\ne c$, a positive weight $w_{e,a}$.
Write
\[
 \defin{\omega(u)}\coloneqq\prod_{e:u_e\ne c}w_{e,u_e},\qquad
 Q_c(\mathbf x;\mathbf w)
 =\sum_{u\in\mathcal F}\omega(u)\prod_{e\in B(u,c)}x_e.
\]
This is the \href{rank.htm#completionPolynomials}{weighted completion
polynomial}. For a set $T\subseteq E$ of size $r_c$, let $q_T$ denote
the coefficient of $\prod_{e\in T}x_e$ in $Q_c$.

Recall that a nonzero real polynomial is
\href{series-parallel.htm#seriesParallelLift}{real stable} if it does
not vanish when every variable has positive imaginary part.

\begin{definition}[Relatively simple vertex]
A vertex of a polytope of affine dimension $d$ is
\defin{relatively simple} if exactly $d$ edges contain it.
We call $u\in\mathcal F$ a \defin{simple root} when its assignment vector
is relatively simple in $P(\mathcal F)$.
\end{definition}

The dimension is relative to the affine hull, not the full space of cells.
The \href{polytope.htm#polytopeEdges}{cycle-edge theorem} identifies the
incident edges with the feasible cyclic exchanges at $u$.

\begin{theorem}[Simple-vertex stability and matching representation]
Suppose that $\mathcal F$ has a simple root $u$. Put
$B=B(u,c)$ and $A=E\setminus B$. Then $q_B>0$, and
\[
 \begin{gathered}
 \defin{W_{ba}}\coloneqq\frac{q_{(B\setminus\{b\})\cup\{a\}}}{q_B}
 \qquad(b\in B,\ a\in A),\\
 Q_c(\mathbf x;\mathbf w)
 =q_B\sum_M\left(\prod_{(b,a)\in M}W_{ba}\right)
       \prod_{e\in(B\setminus B(M))\cup A(M)}x_e.
 \end{gathered}
\]
Here $M$ ranges over bipartite matchings between $B$ and $A$, and
$B(M),A(M)$ are their endpoint sets. An entry $W_{ba}$ is zero if that
single replacement has no completion. In particular, $Q_c$ is real stable.
\end{theorem}

\begin{proof}[Sketch]
At a simple root the incident cycle directions are linearly independent.
Root coordinates express each word uniquely as a position-disjoint
selection of feasible root exchanges. The exchanges touching the target
color remove one position from $B$ and add one from $A$; their compatible
selections form bipartite matchings.

The remaining exchanges use hidden colors only. Their connected conflict
components have clique boundaries among the target exchanges. Thus a
matching meets each such boundary at most once. Summing the hidden
choices component by component gives a positive common factor and a
nonnegative activity for each target edge. The empty matching has
coefficient $q_B$, and the one-edge matchings give the displayed ratios.

To see stability from the representation, introduce a hidden color $h_a$
of quota one for each $a\in A$. At $b\in B$, allow $c$ and the colors
$h_a$ with $W_{ba}>0$, giving the latter weight $W_{ba}$. At $a\in A$,
allow only $c,h_a$, with hidden weight one. Give $c$ quota $|B|$.
These assignments are in weight-preserving bijection with the matchings:
$h_a$ moves from $a$ to $b$ precisely when $(b,a)\in M$.
The \href{series-parallel.htm#listStability}{allowed-list stability
proposition} proves stability of this completion polynomial. Multiplying
by $q_B>0$ proves the claim.
\end{proof}

\begin{remark}[Formalization status]
% Related Lean: ColoredMatroids/SimpleVertex/Stability.lean,
% marginal_isStable_ne_zero_of_simple and marginal_eq_matching_fullList.
E.~Barkeling reports a Lean formalization of the simple-vertex
criterion and matching representation. The supplied exposition marks its
statement review against the formalization as pending. The corollaries
below are paper deductions; no formalization of those deductions is
asserted here.
\end{remark}

Consequently the weighted completion polynomial is Lorentzian
\cite[Proposition~2.2]{BrandenHuh2020}. This proves the
\href{rank.htm#lorentzianConjecture}{completion-polynomial conjecture}
for families with one simple assignment vertex, with no bound on the
number of colors or their quotas. For nonnegative hidden weights, the
same conclusions follow by limits whenever the polynomial is nonzero.

\section[stableIntersectionCounts]{Real-rooted intersection counts}

Fix $S\subseteq E$. We count how many target positions lie in $S$, rather
than only which target sets occur.

\begin{corollary}[Intersection-count polynomials]
Under the simple-root hypothesis and positive hidden weights, the polynomial
\[
 \defin{F_S(t)}\coloneqq
 \sum_{v\in\mathcal F}\omega(v)t^{|B(v,c)\cap S|}
\]
is nonzero and has only real nonpositive zeros.
\end{corollary}

\begin{proof}
Set $x_e=t$ for $e\in S$ and $x_e=1$ otherwise. Specialization at a real
constant and identification of variables preserve stability or give zero.
Positive weights exclude zero. A real univariate stable polynomial has
only real zeros, and nonnegative coefficients exclude positive zeros.
\end{proof}

In particular, if $d=\deg F_S$, its coefficients are ultra-log-concave
of order $d$ with no internal zeros, by Newton's inequalities.
After normalization, $F_S(t)/F_S(1)$ is a product of factors
$1-p+pt$, $0\leq p\leq1$. Thus the intersection count has the same
distribution as a sum of independent Bernoulli variables; the individual
target-position indicators need not be independent.

For the complement $A$ of the root target set, the representation gives
\[
 F_A(t)=q_B\sum_M\left(\prod_{(b,a)\in M}W_{ba}\right)t^{|M|}.
\]
The number of target positions moved away from the root is therefore
enumerated by a weighted matching polynomial.

\section[rookPolynomialCorollary]{The classical matching and rook corollary}

We can also go in the other direction: start with a bipartite graph and
recover its matching polynomial from a simple assignment vertex.

Let $G$ have parts $L,R$, with a nonnegative edge weight $a_{\ell r}$.
Define its \defin{matching generating polynomial} by
\[
 \defin{R_G(t)}\coloneqq
 \sum_{M\text{ matching in }G}
       \left(\prod_{\ell r\in M}a_{\ell r}\right)t^{|M|}.
\]
This is the size-generating convention, rather than the alternating
matching-defect polynomial.

\begin{corollary}[Weighted bipartite matching polynomials]
Every zero of $R_G$ is real and strictly negative.
\end{corollary}

\begin{proof}
First take positive weights. Use the
\href{stable-sets.htm#matchingPolytopes}{bipartite matching construction}
with positions $L\sqcup R$, target color $c$ of quota $|L|$, and a hidden
color $h_r$ of quota one for every $r\in R$. At $\ell\in L$, allow $c$
and the colors $h_r$ of its neighbors; at $r\in R$, allow $c,h_r$.
Give cell $(\ell,h_r)$ weight $a_{\ell r}$ and $(r,h_r)$ weight one.

The root word has target color on all of $L$ and color $h_r$ at $r$.
It corresponds to the empty matching. In edge coordinates the matching
polytope has full dimension $|E(G)|$, since it contains zero and every
single-edge indicator. Its tangent cone at zero is the nonnegative
orthant: all its points have nonnegative coordinates, and every
coordinate ray occurs. Thus zero has exactly $|E(G)|$ incident edges
and is simple. The affine matching construction transfers this property
to the assignment root.

For the word corresponding to $M$, the target class is
\[
 (L\setminus L(M))\cup R(M),
\]
and its hidden weight is $\prod_{\ell r\in M}a_{\ell r}$.
Set $x_\ell=1$ for $\ell\in L$ and $x_r=t$ for $r\in R$.
The intersection-count corollary gives real nonpositive zeros of $R_G$.
Its constant term is one, so zero is not a root. Edges of weight zero
can be deleted, proving the nonnegative case.
\end{proof}

A board is a bipartite graph whose two vertex classes are its rows and
columns, with one edge for each allowed square. Matchings are precisely
nonattacking rook placements, so the corollary recovers real-rootedness
of weighted rook polynomials; see A.~Nijenhuis \cite{Nijenhuis1976}.

\begin{example}[The two-by-two board]
For the full two-by-two board with all weights one, there is one empty
placement, four one-rook placements and two two-rook placements. Hence
\[
 R_{K_{2,2}}(t)=1+4t+2t^2,
\]
whose zeros are $-1\pm1/\sqrt2$.
Its matching polytope has dimension four and a simple empty-matching
vertex, although it has no permutation-matroid realization by the
\href{stable-sets.htm#evenHoleObstruction}{even-hole obstruction}.
Repeated color quotas are essential to this realization.
\end{example}

\begin{remark}[Relation to the classical theorem]
This recovers the bipartite case of the classical matching theorem,
not the arbitrary-graph Heilmann--Lieb theorem. The proof of the
simple-vertex criterion already uses allowed-list stability, which
also proves this matching specialization directly. The geometric
criterion explains a common structure; it is not an independent new
proof of the classical result.
\end{remark}
